% Methods — DESI paper v3 (complete rewrite)
% Fixed: PCHMM→PWL (composite likelihood, not HMM); added parameter scan

\subsection*{Data}

We used phased whole-genome sequence data from the 1000 Genomes Project high-coverage dataset~\citep{Byrska2022}, comprising 3,202 individuals from 26 sub-populations. We selected 240 individuals (48 per continental super-population: AFR, EAS, EUR, SAS, AMR) by random sampling from each super-population, balanced to represent constituent sub-populations. Variant calls on all 22 autosomes were obtained from the Project's public release (\texttt{NYGC30x} VCF files, GRCh38/hg38). We restricted analysis to biallelic SNPs with FILTER$=$PASS, no missing genotypes in our 240 samples, and minor allele count $\geq 1$.

\subsection*{Pairwise Windowed Likelihood (PWL)}

DESI estimates within-group mean pairwise TMRCA from per-window aggregate heterozygosity using the Pairwise Windowed Likelihood (PWL), a composite likelihood approach. This method is \textit{not} a hidden Markov model and does not model Markov transitions between genomic windows; each window is treated independently.

\paragraph{Heterozygosity aggregation.}
For each non-overlapping 100 kb window $w$ and each super-population group $g$, we count the total number of heterozygous sites across all ordered pairs of haplotypes:
\[
n_w^{(g)} = \sum_{i < j} \mathbb{1}[\text{haplotypes } i,j \text{ differ at a biallelic SNP in window } w]
\]
where the sum is over all $\binom{2|g|}{2}$ ordered haplotype pairs. The number of such pairs is $M = \binom{2 \times 48}{2} = 2{,}256$ for a 48-individual super-population. Each window contributes a single count $n_w^{(g)}$ and a callable-site count $L_w$ (the number of positions with no missing genotypes in the window, used as the window length denominator).

\paragraph{Poisson likelihood.}
Under the infinite-sites coalescent, the expected count for a group with mean pairwise TMRCA $T$ ka is:
\[
\mathbb{E}[n_w^{(g)}] = 2\,\mu\,L_w\,T \cdot \frac{1{,}000}{G}
\]
where $\mu = 1.2 \times 10^{-8}$ per bp per generation and $G = 28$ years per generation. Treating $n_w^{(g)}$ as Poisson-distributed with mean $\lambda_w = 2\,\mu\,L_w\,T / G \cdot 1{,}000$, the maximum-likelihood estimate for window $w$ is:
\[
\hat{T}_w^{(g)} = \frac{n_w^{(g)} \cdot G}{2\,\mu\,L_w \cdot 1{,}000}
\]
The genome-wide mean TMRCA for group $g$ is $\bar{T}^{(g)} = \text{mean}_w(\hat{T}_w^{(g)})$ over all valid windows (minimum $L_w \geq 50{,}000$ bp). Uncertainty is quantified using a block jackknife with 1 Mb non-overlapping blocks~\citep{Kunsch1989}: the jackknife standard error is $\text{SE} = \sqrt{(K-1)/K \sum_{k=1}^K (\bar{T}^{(g)}_{-k} - \bar{T}^{(g)})^2}$ where $K$ is the number of 1 Mb blocks.

\paragraph{Callable-site correction.}
We apply a callable-site correction using sample-level missingness masks. For each window and each individual, we count the number of non-missing positions from the VCF's FORMAT/GT field; $L_w$ is set to the minimum callable length across all pairs in the group to ensure conservative estimates.

\subsection*{Coalescent simulation and PWL calibration}

We validated PWL using \texttt{msprime}~\citep{Kelleher2016} coalescent simulations. Three demographic models were simulated: Model OoA (standard Out-of-Africa: ancestral $N_e = 15{,}000$, African $N_e = 15{,}000$, OoA bottleneck at 65 ka reducing to $N_e = 3{,}000$), Model A (panmictic: constant $N_e = 10{,}000$ with random group labels), and Model D (OoA structure with reduced ancestral $N_e = 12{,}000$). Each model was simulated three times (seeds 42, 123, 456) with sequence length $10^8$ bp and sample size 48 individuals per group. Exact TMRCA values were obtained from \texttt{tskit} tree-sequence statistics.

\subsection*{Bottleneck parameter scan}

To directly test FitCoal's bottleneck parameterization~\citep{Hu2023}, we performed a parameter scan across a grid of bottleneck severity and duration. Demographic models were implemented using \texttt{msprime.Demography} with three populations: AFR (modern $N_e = 15{,}000$), EAS (modern $N_e = 3{,}000$ post-OoA bottleneck at 65 ka), and ANC (ancestral population). All three modern populations merge into ANC at 65 ka (OoA split time). The bottleneck is applied to ANC between times $(813 + d)$ ka and $813$ ka, where $d$ is the duration in ka.

We scanned:
\begin{itemize}
  \item $N_e^{\text{bn}} \in \{300, 500, 800, 1{,}280, 2{,}000, 4{,}000, 8{,}000, 20{,}000\}$ (8 values, approximately log-spaced)
  \item $d \in \{30, 50, 80, 117, 150, 200, 300\}$ ka (7 values, linearly spaced)
  \item Ancestral $N_e \in \{20{,}000, 50{,}000\}$ (2 values)
\end{itemize}

This yields 112 parameter combinations. Each scenario was simulated 4 times (seeds 0--3) with sequence length $1.5 \times 10^9$ bp (yielding 150 windows of 100 kb each per simulation) and 10 haplotypes per population ($n = 5$ individuals, sufficient to estimate mean TMRCA). Per-window mean TMRCA was computed using \texttt{tskit.diversity(mode="branch", windows=...)} for both AFR and EAS sample sets, then converted to years: $\bar{T}_\text{ka} = d_\text{branch}/2 \cdot G / 1{,}000$, where $d_\text{branch}$ is the branch-mode pairwise diversity (equal to $2 \cdot \mathbb{E}[T_\text{MRCA}]$ in generation units). Two summary statistics were computed for each scenario: (1) the fraction of windows with $\bar{T}_\text{AFR} > 930$ ka, denoted $P(\bar{T}_\text{AFR} > 930)$; and (2) the mean within-AFR TMRCA $\bar{T}_\text{AFR}$. The four seeds were pooled, giving 600 windows per scenario. The FitCoal exact point corresponds to $N_e^{\text{bn}} = 1{,}280$, $d = 117$ ka, ancestral $N_e = 50{,}000$.

Statistical comparison used a Z-test: $Z = (P_\text{obs} - P_\text{sim})/\sigma_\text{sim}$, where $P_\text{obs} = 0.707$ (observed genome-wide), $P_\text{sim} = 0.615$ (FitCoal scenario), and $\sigma_\text{sim}$ is the standard deviation of $P$ across the four seeds.

\subsection*{BGS robustness analysis}

We identified putatively neutral windows by excluding all 100 kb windows overlapping any hg38 RefGene annotation $\pm$50 kb. This retained 30.7\% of autosomal windows (8,451/27,507). Within-group TMRCA was re-estimated in this subset using the same PWL procedure.

\subsection*{Population samples and sub-population analyses}

For AMR sub-population analyses, we used all available ACB (African-Caribbean, $n = 96$), PEL (Peruvian, $n = 85$), and PUR (Puerto Rican, $n = 104$) individuals from the 1000 Genomes Project. Within-group TMRCAs were estimated as above for each sub-population separately. Mutation rate sensitivity analysis was performed by re-running PWL with $\mu \in \{1.0, 1.2, 1.4, 1.6\} \times 10^{-8}$ on chromosome 21 and scaling genome-wide results accordingly.

\subsection*{Software and code availability}

All analyses were performed in Python 3.10 using \texttt{msprime} v1.2~\citep{Kelleher2016}, \texttt{tskit} v0.5, \texttt{numpy}, \texttt{scipy}, and \texttt{matplotlib}. All analysis code is available at \url{https://github.com/EvoClaw/DESI}.
